% ============================================================
% 02_model.tex — Mô hình hoá bài toán AI
% ============================================================

\section{Mô hình hoá Bài toán AI}
\label{sec:model}

Theo yêu cầu bài tập lớn, hệ thống được mô hình hóa qua 4 khía cạnh AI: \textbf{State} (biểu diễn), \textbf{Knowledge} (tri thức), \textbf{Uncertainty} (không chắc chắn), và \textbf{Learning} (học máy).

\subsection{Biểu diễn Trạng thái (State Representation)}

Bài toán N-Days Meal Planning được mô hình hóa:

\begin{itemize}
  \item \textbf{State} $S = (\text{day}, C_{\text{total}}, \text{Cal}_{\text{total}}, R_{\text{selected}})$: Bộ bốn gồm ngày hiện tại, tổng chi phí, tổng calo, và \textit{sorted multiset} các công thức đã chọn.

  \item \textbf{Initial State} $S_0 = (0, 0, 0, \emptyset)$: Chưa chọn công thức nào.

  \item \textbf{Action} $a_k$: Chọn công thức $r_k$ cho ngày tiếp theo:
  \begin{equation}
    S' = (\text{day}+1,\; C + \text{cost}(r_k) + \text{penalty}(r_k),\; \text{Cal} + \text{cal}(r_k),\; R \cup \{r_k\})
  \end{equation}

  \item \textbf{Goal Test} $G(S)$:
  \begin{equation}
    \text{day} = N \;\wedge\; C_{\text{total}} \leq B \;\wedge\; \text{Cal}_{\min} \leq \text{Cal}_{\text{total}} \leq \text{Cal}_{\max}
  \end{equation}

  \item \textbf{Path Cost} $g(S) = \sum_{r \in R_{\text{selected}}} \text{cost}(r) + \text{penalty}(r)$

  \item \textbf{Canonical Form:} $R_{\text{selected}}$ được sort trước khi nạp vào visited set, gộp các trạng thái hoán vị thành một.
\end{itemize}

\subsection{Biểu diễn Tri thức (Knowledge Representation)}

Tri thức hệ thống được biểu diễn bằng \textbf{Horn Clauses} (logic mệnh đề), lưu trong \texttt{data/rules.json}:

\textbf{Quy tắc sức khỏe} (7 rules):
\begin{equation}
  \text{IF condition:}h \in \text{WorkingMemory} \Rightarrow \text{exclude\_ingredient:}i_1, \ldots, \text{exclude\_tag:}t_1, \ldots
\end{equation}

Ví dụ cụ thể (tiếng Anh, khớp với Kaggle dataset):
\begin{lstlisting}[language=Python, caption=Horn Clause trong rules.json]
{
  "rule_id": "R_STOMACH",
  "antecedents": ["condition:stomachache"],
  "consequents": ["exclude_ingredient:chili",
                   "exclude_ingredient:garlic",
                   "exclude_tag:spicy",
                   "warning:Avoid spicy food"]
}
\end{lstlisting}

\subsection{Mô hình hoá Không chắc chắn (Uncertainty)}

Yếu tố không chắc chắn được xử lý bằng \textbf{Extended Noisy-OR Bayesian Network}:

\begin{enumerate}
  \item \textbf{Posterior Preference} --- $P(\text{like} \mid F)$: Dự đoán sở thích người dùng dựa trên features (budget fit, health match, cuisine match, tag overlap).

  \item \textbf{Digestive Risk} --- $P(\text{risk} \mid I, H)$: Đánh giá rủi ro tiêu hóa dựa trên nguyên liệu $I$ và tình trạng sức khỏe $H$, sử dụng Noisy-OR model với Synergy Hidden Causes.
\end{enumerate}

\begin{equation}
  P(\text{risk}) = 1 - \prod_{i=1}^{k}(1 - P_i) \times \prod_{j=1}^{m}(1 - P_{\text{synergy}_j})
\end{equation}

Trong đó $P_i$ là xác suất rủi ro độc lập của yếu tố $i$, và $P_{\text{synergy}_j}$ là xác suất rủi ro của tổ hợp hiệp đồng $j$ (hidden cause). Kết quả luôn $\in [0, 1]$ mà không cần clamping.

\subsection{Học máy (Machine Learning)}

Hai thành phần ML chính:

\begin{enumerate}
  \item \textbf{NER + Feature Extraction (TF-IDF):}
    \begin{itemize}
      \item \textbf{Input:} Văn bản tiếng Anh tự do
      \item \textbf{Output:} Vector thực thể $E = (e_{\text{budget}}, e_{\text{health}}, e_{\text{ingredients}}, e_{\text{cuisine}})$
      \item Sử dụng Ontology-based Entity Extraction + regex patterns
      \item TF-IDF vectorization cho recipe matching (Cosine Similarity)
    \end{itemize}

  \item \textbf{Multinomial Naive Bayes Classifier (Cuisine Classification):}
    \begin{itemize}
      \item \textbf{Bài toán:} $P(C \mid I_1, \ldots, I_n)$ --- phân loại cuisine dựa trên danh sách nguyên liệu
      \item \textbf{Training data:} Kaggle Food.com dataset (hàng nghìn recipes qua stratified sampling) hoặc \texttt{recipes.json} (25 recipes) khi chạy local
      \item Laplace smoothing + log-sum-exp normalization
    \end{itemize}
\end{enumerate}

\subsection{Tổng hợp: Ánh xạ 5 Thành phần Bắt buộc}

\begin{table}[H]
\centering
\begin{tabular}{|c|l|l|l|}
\hline
\textbf{Ký hiệu} & \textbf{Yêu cầu} & \textbf{Module} & \textbf{Thuật toán} \\
\hline
(A) & Biểu diễn \& Tìm kiếm & \texttt{search\_engine.py} & A* Search \\
\hline
(B) & Heuristic / CSP & \texttt{csp\_solver.py} & Backtracking + MRV + FC \\
\hline
(C) & Logic / Tri thức & \texttt{knowledge\_base.py} & Forward Chaining \\
\hline
(D) & Bayes / Xác suất & \texttt{bayes\_risk.py} & Extended Noisy-OR \\
\hline
(E) & Học máy & \texttt{ml\_classifier.py}, \texttt{feature\_extractor.py} & Naive Bayes + TF-IDF \\
\hline
\end{tabular}
\caption{Ánh xạ 5 thành phần bắt buộc $\to$ modules dự án}
\label{tab:5pillars}
\end{table}
